\documentclass[11pt,letterpaper]{article}
\usepackage[margin=0.92in,headheight=15pt]{geometry}
\usepackage[T1]{fontenc}
\usepackage{lmodern,microtype,amsmath,amssymb,amsthm}
\usepackage{booktabs,longtable,tabularx,array,xcolor,enumitem,listings}
\usepackage{xurl,hyperref,fancyhdr}
\definecolor{ink}{RGB}{27,53,77}
\definecolor{accent}{RGB}{19,99,122}
\hypersetup{colorlinks=true,linkcolor=ink,urlcolor=accent,citecolor=ink,
 pdftitle={U208: a three-cell fold for the consistency of PA},
 pdfauthor={Research prepared for Vladimir Reshetnikov}}
\pagestyle{fancy}\fancyhf{}\lhead{\small U208: a three-cell fold}\rhead{\small 29 October 2026}\cfoot{\thepage}
\setlength{\parindent}{0pt}\setlength{\parskip}{5pt}
\setlist{itemsep=3pt,topsep=4pt}
\newcommand{\PA}{\mathsf{PA}}
\newcommand{\U}{\mathsf U}
\newcommand{\W}{\mathsf W}
\newcommand{\Con}{\operatorname{Con}}
\newcommand{\Halt}{\operatorname{Halt}}
\newcommand{\Tape}{\mathcal T}
\newcommand{\act}[3]{#1\,#2\,#3}
\newcommand{\code}[1]{\texttt{#1}}
\newtheorem{theorem}{Theorem}[section]
\newtheorem{lemma}[theorem]{Lemma}
\newtheorem{corollary}[theorem]{Corollary}
\lstset{basicstyle=\ttfamily\small,breaklines=true,columns=fullflexible,
 keepspaces=true,frame=single,framerule=0.3pt,rulecolor=\color{black!25}}
\begin{document}
\begin{center}
{\color{ink}\LARGE\bfseries U208: a three-cell fold\\[5pt]
for the consistency of Peano arithmetic}\par
\vspace{10pt}
{\large Five U209 proposals, two machines, one further state deletion}\par
\vspace{8pt}
Research prepared for Vladimir Reshetnikov\quad\textbullet\quad29 October 2026
\end{center}

\begin{abstract}
The five code-bearing 209-state proposals in the specified ProveIt snapshot
represent two rooted, action-labeled transition graphs. Four are the same
machine up to state renaming; the Compatibility package contains the other
machine. We compare their source layouts, proof scopes, reconstruction
programs, and freshly reproduced verification results.

Starting from the Joint/Compatible-Folds table, we delete native state 211
and reuse state 172. Only three retained transition rows change. The proof
uses information from the preceding execution: state 211 scanning zero has
a one immediately to its left, whereas state 172 has the predecessor word
$10$. Four finite paths inspect at most the current cell and its two left
neighbors, preserve the whole tape, and simulate each old step by one, three,
or seven new steps. This establishes halting equivalence on every common
initial tape from state 0, independently of the source program and theory
calibration.

The resulting complete binary, single-tape machine has \textbf{208 working
states and 416 transitions}, including initialization. All 5,084
arbitrary-value primitive contracts were freshly rederived from its literal
table, covering every transition row; the complete 48,140-step blank startup
is unchanged. The new finite reduction is proved here. The proposed
$\Con(\U)\leftrightarrow\Con(\PA)$ endpoint retains the incoming mathematical
calibration and is not a newly kernel-checked PA, Lean, Rocq, or MM0 theorem.
\end{abstract}

\newpage
\begingroup\small\setlength{\parskip}{0pt}
\tableofcontents
\endgroup
\newpage
\section{Result, conventions, and proof boundary}

The main deliverable is \path{machine/u208.tm}; Appendix~\ref{app:table}
prints its complete table. It is obtained from the Joint U209 table by a
deterministic edit and a checked dense renaming.

\begin{center}
\begin{tabular}{@{}ll@{}}
\toprule
Working states & $0,1,\ldots,207$\\
Separate halt label & $208$\\
Alphabet and blank & $\{0,1\}$, blank $0$\\
Tape & One two-way infinite tape\\
Initial configuration & State $0$, head at $0$, all cells $0$\\
Transition format & $q\ b\ w\ D\ r$: write, move, enter $r$\\
Movement & Every row moves exactly one cell, left or right\\
Initialization & Included in the 208 working states\\
Total rows & $416$, with no missing read case\\
\bottomrule
\end{tabular}
\end{center}

Entering the separate halt label halts after the indicated write and move.
The companion \path{machine/u208.native.tm} retains sparse historical state
names and uses 254 as its halt label. A native name such as 211 is an
identifier, not a state count. All state numbers used in the proof below are
\emph{Joint native labels}. The dense table is the submission-ready object.

Let $P$ denote the exact Joint U209 parent and $M$ the new U208 machine.
The new result is
\begin{equation}\label{eq:new}
 \PA\vdash\Halt(M,\mathrm{blank})\leftrightarrow
                  \Halt(P,\mathrm{blank}).
\end{equation}
More generally, the ordinary mathematical proof applies to every common
initial tape from state 0. For first-order arithmetization one uses coded
finite-support tapes and finite run histories; the blank-tape instance in
\eqref{eq:new} is immediate. No consistency assumption is used.

The supplied proposals identify the parent's nonhalting with consistency by
their retained numerical-calculus, enumeration, and interpretation arguments.
Composing that calibration with \eqref{eq:new} gives the intended endpoint
\begin{equation}\label{eq:endpoint}
 \PA\vdash\neg\Halt(M,\mathrm{blank})
 \leftrightarrow\Con(\U)\leftrightarrow\Con(\PA).
\end{equation}
The new fold does not strengthen the verification status of that inherited
calibration. This report supplies a complete local reduction proof and
executable checks. It does not supply a new end-to-end proof-assistant
certificate, prove consistency, or infer nonhalting from a finite run.

\begin{center}\small
\begin{tabularx}{\linewidth}{@{}>{\raggedright\arraybackslash}p{36mm}X@{}}
\toprule
Claim & Evidence and scope\\
\midrule
Exact 208-state object & Complete native/dense tables, totality, target
closure, graph reachability, bijective renaming, deterministic reconstruction.\\
New $P\leftrightarrow M$ reduction & Complete incoming-history analysis,
four local paths, 414 ordinary-row identities, positive-block simulation.\\
Literal source implementation & Fresh universal symbolic replay of all
5,084 contracts and exact whole-tape blank initialization.\\
Consistency endpoint & Composition with the incoming source/calculus and
U/PA proof translations; not newly formalized here.\\
\bottomrule
\end{tabularx}
\end{center}

\section{What the supplied proposals contain}
\subsection{Pinned inventory}

The comparison is frozen at ProveIt commit
\begin{center}\small
\path{fe8b3a2d1db0033113e986590f79ecba3e5ef500}.
\end{center}
All eight entries in the downloaded directory inventory, including its
README, were checked against their Git blob identities. The seven ZIPs are
the five code-bearing U209 proposals, the extra Joint PDF-only archive, and
\path{pa_unary_power_cover.zip}. Exact sizes, download URLs, Git object
identities, and SHA-256 hashes are retained in \path{provenance/}.
The five input construction reports are references~\cite{plain,composed,folds,joint,compat}.

\begin{center}\small
\begin{tabularx}{\linewidth}{@{}p{28mm}cp{16mm}X@{}}
\toprule
Proposal & Class & Depth & Distinctive presentation\\
\midrule
Plain U209 & A & 42 & Thirty rotations from U212; an additional obstruction
to removing a dominated propositional axiom row.\\
Composed & A & 42 & Thirty rotations with a label permutation; continuation
transport and calibration exposition.\\
Compatible Folds & A & 42 & Staged emitter, combined, and prefold tables;
separate local and source audits.\\
Joint & A & 42 & Direct allocation of an 18-selector region; compact
deterministic recipe and complete retained lineage.\\
Compatibility & B & 41 & A different allocation across a shared helper;
Research U211 to U210 to U209.\\
\bottomrule
\end{tabularx}
\end{center}

The four Class A packages encode the \emph{same literal machine} after
renaming. Compatible Folds and Joint have byte-identical native tables,
dense tables, and final source maps. Plain and Composed use other native
names for a few selectors. The extra
\path{U209_ConPA_Joint_2026-10-29 (1).zip} contains only a PDF, which identifies
the Plain table by its hashes; it is not a sixth implementation.

\subsection{How equality and difference were established}

The independent comparison program performs a breadth-first traversal from
state 0, always visiting the read-zero edge before the read-one edge. It
assigns canonical names on first encounter and preserves the read, write,
move, and halt information. All 209 working states are reachable in each
transition graph. This is graph reachability, not a claim that every state
is visited on the blank computation.

For each of the ten pairs of packages, the checker also constructs the
forced state bijection and verifies every transition equation. Four
packages have the same canonical graph; attempted bijections to
Compatibility fail with a recorded finite conflict. This establishes two
isomorphism classes rather than merely two groups of file hashes.

For example, the nontrivial Joint-to-Plain native renaming is
\[
129\mapsto217,\quad204\mapsto215,\quad214\mapsto129,\quad
216\mapsto204,\quad217\mapsto214,\quad215\mapsto216.
\]
Every other label is fixed. Consequently the new helper 217 in this report
is called 214 in Plain. Copying the numeric patch across packages without
this renaming would produce a different, unjustified machine.

All five source programs have the same ordered 3,909 primitives and operands
and the same 5,083 guarded ordinal control edges. Their source maps have
197 distinct syntax subtrees and maximum zero-run length six. Classes A
and B differ in precisely 720 instruction addresses, the ordinal interval
$[3016,3736)$. Equality of the ordinal graph was checked independently by
binary-address increment, ancestor stripping, and zero-padded target
decoding, including the absence of stray ones beyond the decoded leaf.

\subsection{The shared earlier savings}

The proposals all inherit three savings from the same U212 lineage. Class A
deletes native states $180,181,218$ by the identifications
\[
 180\mapsto163,\qquad181\mapsto220,\qquad218\mapsto157.
\]
The first two folds use a bit written by the incoming transition to
distinguish two old roles. The third substitutes an existing three-step
behavior for a decoder entered with a predecessor zero. Reassociation is
needed because the separately proposed U211 patches do not coexist on their
original layouts. In Joint labels the retained selector rows include
$163/0\mapsto\act0L{211}$ and $220/0\mapsto\act0L{129}$.

Compatibility reaches a different arrangement. Its final contextual fold
is $180\mapsto216$, and the corresponding changed row is
$216/0\mapsto\act0L{211}$. A backward call into helper 129 reads a known
zero on the return leg. That call uses only rows $158/0$ and $204/0$ after
the tag test, so row $204/1$ remains available for changing the source layout.
The eight-factor reassociation exploits exactly this restricted interface.

These observations distinguish useful mechanisms from presentation changes.
Renaming Class A or giving another sequence of rotations to it cannot
reduce its state count. A new deletion is needed. The present construction
deletes the already shared helper 211 itself and compensates for its two
entry directions using a stronger execution-history invariant.

\subsection{Fresh reproduction and runtime comparison}

All five supplied quick reproduction commands passed in isolated working
copies. One representative of each isomorphism class was then replayed
fully. No saved success flag was substituted for these executions.

\begin{center}\small
\begin{tabular}{@{}lrr@{}}
\toprule
Fresh full verification & Joint, Class A & Compatibility, Class B\\
\midrule
Arbitrary-value primitive contracts & 5,084 & 5,084\\
Literal rows covered & 418/418 & 418/418\\
Symbolic inference applications & 1,250,049 & 1,248,748\\
Exact blank startup steps & 48,140 & 48,140\\
Fresh cases equal packaged records & Yes & Yes\\
\bottomrule
\end{tabular}
\end{center}

The two classes do not have a uniform runtime ordering. Their per-primitive
affine duration expressions have equal nonconstant coefficients, but
different constants. For the same logical registers and gap/spare
parameters, Compatibility is faster on 329 contracts, slower on 466, and
equal on 4,289. The differences range from $-6$ to $+70$ literal steps.
The histogram and case comparison are retained in
\path{evidence/duration_comparison.json}. A smaller maximum source depth or
fewer symbolic inference applications is therefore not an overall speed
claim. The actual frequency of those primitive cases is not determined by
this comparison.

Joint is selected as the main baseline because it is byte-identical to
Compatible Folds, its deterministic construction is compact, and it retains
the complete nested calibration material. This is an auditability choice,
not a statement that Class A is faster than Class B.

\section{The new deletion: three changed rows}
\subsection{Literal patch}

Define the projection
\[
 \rho(211)=172,\qquad \rho(q)=q\quad(q\ne211).
\]
Delete the two rows whose source state is 211. Redirect every retained
target through $\rho$. Finally replace the read-zero row at 172. The entire
change to the retained native rows is

\begin{center}
\begin{tabular}{@{}rcll@{}}
\toprule
State & Read & Parent action & U208 action\\
\midrule
163 & 0 & $\act0L{211}$ & $\act0L{172}$\\
163 & 1 & $\act1R{211}$ & $\act1R{172}$\\
172 & 0 & $\act0L{67}$ & $\act0L{217}$\\
\bottomrule
\end{tabular}
\end{center}

The removed rows are $211/0\mapsto\act0R{157}$ and
$211/1\mapsto\act1R{184}$. The retained row
$172/1\mapsto\act1R{184}$ already agrees with the latter. No source
instruction, source address, opcode, axiom formula, numeric convention,
initialization routine, or tape alphabet is changed.

The final table has 208 working states. Its dense-state map is the increasing
bijection from the retained native names to $0,\ldots,207$, extended by
$254\mapsto208$ at the halt. Both directions of this conjugacy are checked
against every row.

\subsection{Why a history predicate is needed}

State 211 has two entrances from state 163. One moves right and writes one;
the other moves left and writes zero. Thus a test that requires every
entrance to a state to write the same predecessor tag cannot justify this
fold. The useful invariant is conditional on the \emph{currently scanned}
symbol, and its proof looks one step farther into the incoming history.

State 172 also moves in a different direction on its zero branch than state
211. A direct row identification is invalid. The replacement must inspect
the predecessor, and in one case the predecessor's predecessor, before it
can recover the correct old action. The tape already contains the required
information; no new symbol or state is needed.

\section{Complete proof of the new machine equivalence}\label{sec:proof}
\subsection{Incoming inventories}

For a configuration $(q,h,\Tape)$, $h$ is the head coordinate and
$\Tape(i)\in\{0,1\}$ is the bit in cell $i$. The following are the complete
incoming-row inventories in the exact parent table. A source/read entry
lists the row that enters the indicated target.

\begin{center}
\begin{tabular}{@{}rll@{}}
\toprule
Target & Incoming source/read & Action entering target\\
\midrule
211 & $163/0$; $163/1$ & $\act0L{211}$; $\act1R{211}$\\
172 & $116/0$ & $\act0R{172}$\\
116 & $215/1$ & $\act1R{116}$\\
163 & $132/1$; $199/1$; $217/0$ & $\act1R{163}$; $\act1R{163}$; $\act0R{163}$\\
217 & $216/1$ & $\act1R{217}$\\
\bottomrule
\end{tabular}
\end{center}

All these states differ from the initial state 0. The inventories are
computed from all 418 parent rows, not inferred from an observed blank run.

\begin{lemma}[Conditional predecessor at 211]\label{lem:211}
In every parent run starting in state 0, if state 211 scans zero at head
position $h$, then $\Tape(h-1)=1$.
\end{lemma}
\begin{proof}
The last transition is from 163. If it is $163/1$, it writes one at $h-1$
and moves right to $h$, proving the assertion immediately.

Otherwise the last transition is $163/0$, made at $h+1$. It writes zero
there and moves left to $h$. Every entrance to 163 moves right, so the
preceding transition entered 163 from $h$ and wrote the bit now scanned
by 211. That bit is zero. The complete inventory says that the only such
zero-writing entrance is $217/0$. Every entrance to 217 is $216/1$, which
previously wrote one at $h-1$ and moved right to $h$. The intervening two
steps write only at $h$ and $h+1$. The one at $h-1$ is unchanged.

The history cannot terminate prematurely at initialization: 211, 163, and
217 are not state 0. Thus all required preceding transitions exist.
\end{proof}

\begin{lemma}[Two predecessor cells at 172]\label{lem:172}
In every parent run starting in state 0, a visit to state 172 at position
$h$ satisfies $\Tape(h-2)=1$ and $\Tape(h-1)=0$.
\end{lemma}
\begin{proof}
The unique entrance is $116/0$, which writes zero at $h-1$ and moves right.
The unique entrance to 116 is $215/1$, which previously wrote one at
$h-2$ and moved right. The next transition leaves that cell untouched.
Neither 172 nor 116 is the initial state, so the two-step history exists.
\end{proof}

These lemmas concern \emph{parent} configurations at simulation boundaries.
They need not hold at every intermediate visit to the reused state 172
inside a child block. Confusing those two roles would make a circular
invariant argument. The following literal paths handle the intermediate
visits explicitly.

\subsection{Four exhaustive local cases}

The unchanged helper rows used by the new blocks are
\begin{align*}
217/0&\mapsto\act0R{163}, &217/1&\mapsto\act1R{215},\\
215/0&\mapsto\act0R{157}, &157/0&\mapsto\act0L{67}.
\end{align*}
The child has $163/0\mapsto\act0L{172}$ and
$172/0\mapsto\act0L{217}$ after the patch. The shared read-one row is
$172/1\mapsto\act1R{184}$.

Normalize the old head coordinate to zero. The table lists each child path
as $(\text{state},\text{head offset})$, including its endpoint.

\begin{center}\small
\begin{tabularx}{\linewidth}{@{}rrp{24mm}Xr@{}}
\toprule
Old state & Read & Known cells & Child path & Steps\\
\midrule
211 & 0 & $\Tape(-1)=1$ &
$(172,0)\to(217,-1)\to(215,0)\to(157,1)$ & 3\\
211 & 1 & -- & $(172,0)\to(184,1)$ & 1\\
172 & 0 & $\Tape(-2)\Tape(-1)=10$ &
$(172,0)\to(217,-1)\to(163,0)\to(172,-1)\to(217,-2)
\to(215,-1)\to(157,0)\to(67,-1)$ & 7\\
172 & 1 & -- & $(172,0)\to(184,1)$ & 1\\
\bottomrule
\end{tabularx}
\end{center}

The currently scanned bit is the value in the Read column. Every transition
on these paths writes back the bit it reads. The endpoints are exactly the
parent one-step endpoints, with the state projected by $\rho$:
$(157,1)$, $(184,1)$, $(67,-1)$, and $(184,1)$ respectively.
No internal state is the halt. The longest path reads only offsets
$-2,-1,0$. Every other cell is unread and unwritten; this is a whole-tape
identity with an arbitrary exterior, not a test of finitely many chosen
tapes.

The seven-step case revisits state 172 one cell left. At that intermediate
visit the predecessor is one, so the three-step route used for the first
case is followed. Its termination is established by the displayed path;
there is no appeal to an unproved recursive simulation claim.

\subsection{Ordinary steps and positive progress}

For the 414 parent rows whose source state is neither 211 nor 172, the
tables satisfy the exact quotient equation
\begin{equation}\label{eq:ordinary}
\delta_M(\rho(q),b)=(w,D,\rho(r))
\quad\text{whenever}\quad\delta_P(q,b)=(w,D,r).
\end{equation}
The two redirected rows of 163 satisfy this equation. The four exceptional
rows are precisely the four cases above.

\begin{theorem}[All-tape positive-block simulation]\label{thm:fold}
Start $P$ and $M$ in state 0 with the same tape and head coordinate. At
matched boundaries the head and entire tape agree, and the child state is
$\rho$ of the parent state. Every parent step is simulated by a positive
child block of length at most seven. If a parent prefix has $n$ steps, its
matching child prefix has exactly
\begin{equation}\label{eq:clock}
 n'=n+2N_{211,0}(n)+6N_{172,0}(n),
\end{equation}
where $N_{q,b}(n)$ counts parent steps in that prefix executed at state $q$
while reading $b$. In particular, $n\le n'\le7n$. The two machines have the
same halting behavior from every such common initial configuration.
\end{theorem}
\begin{proof}
Initially $\rho(0)=0$. Induct on parent steps. At an ordinary state,
\eqref{eq:ordinary} supplies one child step with the required endpoint.
At state 211 or 172, the two predecessor lemmas supply exactly the bits
required by the applicable local case. That case supplies the matching
positive block and preserves the entire tape. Concatenation proves the
boundary correspondence. Counting the lengths $3,1,7,1$ gives
\eqref{eq:clock} and its bounds.

If the parent halts, the concatenated child blocks reach the same halt
label, which is fixed by $\rho$. For the reverse implication, suppose the
child halts after at most $t$ steps but the parent has not halted within
$t$ steps. Simulating that parent prefix yields a child prefix of at least
$t$ steps containing no halt: ordinary steps cannot halt unless the parent
does, and no exceptional block has an internal halt. Determinism makes
this the actual child prefix, contradicting the assumed child halt.
Thus a child halt within $t$ steps implies a parent halt within $t$ steps.
This establishes both directions using finite prefixes only.
\end{proof}

\begin{corollary}[PA-formalizable reduction]
For blank input, \eqref{eq:new} holds. The finite-horizon implications can
be expressed in PA as
\[
 \Halt_{\le n}(P)\ \Longrightarrow\ \Halt_{\le7n}(M),
 \qquad
 \Halt_{\le n}(M)\ \Longrightarrow\ \Halt_{\le n}(P).
\]
\end{corollary}
\begin{proof}
The tables and all local cases are fixed finite objects. Histories, tape
updates, and runs have primitive-recursive codes. The predecessor lemmas
inspect at most three transitions; the global proof is induction on a
finite coded run length. The same proof therefore gives the stated PA
derivations. This explains formalizability; it is not an exported PA proof
object or a claim that the Python programs are proof-assistant kernels.
\end{proof}

\section{Verification of the delivered U208 table}
\subsection{Two independently implemented finite checks}

\path{code/check_fold.py} checks the complete incoming inventories, the
required helper rows, all four paths, all 414 ordinary quotient equations,
totality, graph reachability, and the one-state deletion. Its conditions are
explicit exceptions and remain active with Python \code{-O}. It rejects
eight damaged constructions, including wrong direction, wrong write,
wrong helper target, and damaged predecessor conditions.

The separate \path{code/check_fold_independent.py} imports neither the
builder nor the discovery checker, source parser, runtime engine, or saved
verdict. It enumerates all bounded incoming histories and solves their
literal read/write constraints. For state 211 scanning zero there are six
length-three state histories: two are inconsistent with the scanned zero,
and all four remaining histories force the left bit to one. The unique
length-two history into 172 forces the predecessor word $10$. The checker
also excludes entries too close to initialization to have the histories
used by the proof. It independently checks the exact child paths and full
table projection.

The negative controls are error-detection checks, not an additional theorem.
The theorem is the invariant and finite-block argument above. The independent
history analysis reduces the risk that a hand-selected predecessor census
silently omitted an entrance.

\subsection{Fresh arbitrary-value runtime derivation}

The source remains byte-identical to Joint's. Nevertheless, the literal
runtime contracts were freshly derived from the final U208 table. The two
inherited symbolic engine modules are copied unchanged and checked against
their frozen SHA-256 identities before execution.

\begin{center}
\begin{tabular}{@{}lr@{}}
\toprule
Fresh U208 verification item & Result\\
\midrule
Working states / transitions & 208 / 416\\
Source primitive occurrences & 3,909\\
Guarded primitive contracts & 5,084\\
Literal transition coverage & 416 / 416\\
Symbolic inference applications & 1,251,865\\
Minimum affine duration constant & 5\\
Duration coefficients nonnegative & All\\
Primary arithmetic registers & Nine arbitrary naturals\\
Gap parameter & Arbitrary natural\\
Spare registers & Arbitrary nonempty finite list\\
\bottomrule
\end{tabular}
\end{center}

Each decrement primitive has separate zero and positive cases. The engine
uses literal steps, induction over same-state unary scans, four rederived
arbitrary-block traversal laws extended by list induction, and explicit
word normalization. It refuses to inspect a possibly empty symbolic run
without a positivity justification. Every case has positive affine duration;
every nonhalting primitive reaches its expected whole-tape final frame, and
the halt primitive reaches the literal halt label. The complete per-case
records are compressed in \path{evidence/runtime_cases.json.gz}.

This is stronger than a bounded numerical test: register values, gap, and
spare-list length are symbolic. It is still executable mathematical
induction implemented in ordinary Python. The inference implementation has
not been checked by a proof-assistant kernel in this release. The direct
fold theorem does not depend on trusting that engine; the replay provides
an additional check that the emitted final table realizes the inherited
primitive interface.

\subsection{Exact initialization and source correspondence}

The parent and child blank runs agree transition by transition for the first
48,140 steps. They finish at native state 84 with the head at $-2$, header
at 54, source ordinal 136, and source address
\texttt{001100110101010}. The logical register tuple is
\[
 (1,0,13,0,0,1,0,0,1),
\]
with 43 spare registers of value one. The entire finite-support tape has
exactly 120 one-cells; all other cells are zero. The new patch is not visited
during this startup. Both the full support and the complete trace digest
are compared, rather than only a few visible registers.

The inherited startup lands inside a protected axiom-builder continuation,
not at a clean abstract proof-search root. Its correctness depends on the
retained source/calculus startup argument. The new fold preserves that
continuation exactly; it does not reinterpret the initial cached value one
as a theorem. This distinction remains part of the logical calibration.

The exact trace SHA-256 is
\begin{center}\small
\path{41e9a165abae389c3b182ebbf8d8fd1c260d3d36d4f7ebb6d60e235bac40f19a}.
\end{center}

\subsection{Runtime cost of the new saving}

Equation~\eqref{eq:clock} is an exact comparison with the chosen U209 parent
on any corresponding run prefix. This is a state saving, with an at-most
sevenfold step overhead. It is not a claimed speed improvement. The bound
can be sharpened for a particular execution by counting its two exceptional
zero-read cases, but no distributional claim about the consistency search
is needed for correctness.

The source addresses and control graph do not change in this reduction,
so its all-tape scope is stronger than the earlier reassociation stages.
It does not retroactively make those earlier source transformations
all-tape simulations.

\section{Logical calibration: what is inherited}
\subsection{Unary U and its PA interpretation}

The proposed theory U is classical equality-free first-order logic over a
nonempty domain with one binary relation, written membership. Define
\[
 x\subseteq a\quad:\Longleftrightarrow\quad
 \forall t\,(t\in x\to t\in a).
\]
For every formula $\varphi(x)$ with no other free variable, its sole
nonlogical schema is
\begin{equation}\label{eq:U}
\forall a\,\exists b\left[
 \forall x\,(x\in b\to\varphi(x))\ \land\
 \forall x\,(x\subseteq a\to(\varphi(x)\to x\in b))\right].
\end{equation}
Bound variables are freshened. The restriction concerns free matrix
parameters; it is not a restriction to one-variable logic. The schema
provides a weak cover and does not assume extensionality or an exact power
filter for every ambient source.

The unary report~\cite{unary} gives interpretations in both directions and
uniform finite-refutation transformations. In the upper interpretation,
membership is binary-bit membership in a PA model. The subset condition
bounds the candidate elements, and formula-specific finite comprehension
is proved by PA induction. This also proves irreflexivity. The proof does
not evaluate a general truth predicate.

The lower interpretation uses numerical representatives modulo
coextensionality, with saturation providing the needed membership
congruence. Its construction order is consequential: local predecessor
facts first give induction without total successors; comparison and least
numerical failures then give bounded numerical comprehension with the
parameters bound inside legal unary matrices. Total successors, numerical
powers, saturated ordered-pair codes, and canonical bounded histories
follow. Addition and multiplication are defined by those histories, and
translated PA induction uses the earlier induction theorem.

Our reading found no concrete illicit free parameter or circular successor
assumption in those examined steps. The package supplies ordinary
mathematical proof templates, not an implemented general translator for
arbitrary PA proofs or a kernel certificate. We retain that scope rather
than presenting finite diagnostics as a new proof of equiconsistency.

\subsection{The machine-facing calculus is a different presentation}

The U209 machine does not implement a unary-free-variable validator. Its
inherited numerical calculus is calibrated to unrestricted weak covering
$\W$, allowing additional matrix parameters, together with global
irreflexivity. Unary U embeds into this theory; the binary-membership
interpretation in PA validates the additional instances and irreflexivity.
These interpretations bridge the consistency presentations.

The native numerical coding uses dual Cantor pairing with constructor tags
$4P$, $4P+1$, $4P+2$, and $4P+3$. Code zero is a truth alias; code one is
diagonal membership used as falsity under irreflexivity. The retained
calibration covers the logical rows, the dominated propositional row,
modus ponens, generalization, the gated covering/vacuity bundle, and the
exceptional small-code and variable-collision cases. Completeness also
requires compiling proofs above arbitrary previous stack contents and
cofinal successful enumeration.

This is \emph{not} the separate positive dyadic theorem stream described in
the unary-theory ZIP. That report's positive stream is another calibrated
interface; its existence does not mean it has been compiled into the
delivered table. The new U208 reduction preserves the actual U209 source
and hence preserves its existing calibration obligations exactly.

The complete older U222 mathematical material is present in the retained
Joint archive through its selected U212, U213, U214, and U218 ancestry.
The included logical reference sources have unbundled imports and should
not be described as a runnable new U208 Rocq project. Original warnings
superseded by the exceptional-code correction must be read together with
that correction, as the incoming reports explain.

\subsection{Relationship to the supplied ZFC verifications}

The m254 project supplies external Coq endpoints for its literal table,
calculus, and ZF/ZFC target~\cite{m254}. ProveIt's K215 development supplies
external Lean and Rocq endpoints for its exact K215 table~\cite{k215}.
The separate k215-pa repository supplies the internal PA theorem in
Metamath Zero~\cite{k215pa}. These are relevant examples of end-to-end
verification architecture, but none certifies the different U208 table or
its PA calibration. Those large projects were read for scope and were not
rebuilt during this investigation.

\section{Rejected reductions and lessons from the search}
\subsection{An attractive sharing transformation that fails}

A structural search suggested reassociating
\[
 ((98,12),150)\longmapsto(98,(12,150)),\qquad
 ((166,12),150)\longmapsto(166,(12,150)).
\]
At the table level these are rotations of $(155,85)$ and $(52,154)$.
Afterward states 85 and 154 have identical rows, apparently saving one
state by sharing the suffix $(12,150)$. The final identical-row quotient is
valid; the preceding rotation is not justified.

The reason is hidden reuse through backward execution. A purely rightward
source traversal sees three occurrences of 155. Full literal prefix
analysis finds five occurrences of 155 and ten bindings to 85, including
five not under those 155 roots. One such path is
\[
76\to152\to159\to118\to155\to85.
\]
It revisits already consumed bits. The exact source prefix
\texttt{0011001000001}, with a preceding one, exposes the failure: the old
machine completes at state 42 and offset 13 without changing the tape;
the rotated machine reaches state 41 at offset 12 and overwrites a one
with zero. The counterexample is retained in
\path{research/rejected_rotation.json}. No machine from that failed route
is included as a candidate.

This is a useful limitation on source-DAG optimization. Literal equality
of two rows after a source edit proves the quotient stage only. One must
also validate all physical entries to every repurposed decoder, including
entries that reread earlier address bits.

\subsection{Shortening the origin detector also fails directly}

Another test attempted to delete native state 34 from the seven-zero origin
detector by redirecting $23/0$ directly to 75. All source-address six-zero
runs are terminal after trimming the suffix, which initially suggested a
six-zero detector might suffice. The actual runtime uses temporary markers,
so a static trimmed-address statistic is insufficient. The shortened table
already fails the first primitive's zero-decrement symbolic contract at
source address \texttt{000000}. This exploratory table is not delivered.

Both failed ideas were screened against the literal runtime before a state
saving was accepted. The delivered 208-state result is solely the
three-cell history fold proved in Section~\ref{sec:proof}.

\section{Further research questions}
\begin{enumerate}
\item \textbf{Can bounded-history certificates become a systematic optimizer?}
The successful predicate couples the current read bit with earlier writes;
it is missed by a uniform one-cell incoming-tag test. Enumerating feasible
short transition histories, solving bit constraints, and synthesizing local
replacement paths provides a finite search space with independently
checkable certificates. The useful next objective is one more sound deletion,
not a claim that the current search has exhausted all 207-state possibilities.

\item \textbf{Can the generic fold theorem be formalized once?}
A proof-assistant library could accept a state projection, finite incoming
history covers, ordinary-row equations, and finite positive blocks, then
produce the blank-halting equivalence for a literal table. This would make
the incremental U209-to-U208 proof small. It would not by itself formalize
the separate numerical-calculus and U/PA interpretation layers.

\item \textbf{Can source sharing account for every physical decoder entry?}
The rejected rotation suggests enriching each source node with its complete
literal prefix interface, including known previously read bits and backward
entries. A transformation should transport those interfaces before counting
a shared suffix as a state saving.

\item \textbf{Do the two decoder allocations offer different next savings?}
Class B has a different 720-primitive address region and maximum depth 41.
The present history-fold mechanism also checks locally in its native
layout, with the appropriate helper names. The two layouts should remain
distinct search inputs rather than being counted as five independent designs.

\item \textbf{Which inherited calibration step should receive the next kernel proof?}
Priority candidates are the numerical-calculus refutation bridge, protected
nonzero startup, and the unary U/PA proof translations. These are precisely
where an incremental table certificate cannot replace an end-to-end proof.
\end{enumerate}

No lower bound or global optimality assertion follows from this investigation.
The achieved result is an explicit, reproducible improvement from 209 to 208
working states with a small new halting-equivalence proof.

\section{Reproduction and archive map}

Use Python 3.10 or later with its standard library and enabled assertions.
The default command is offline and writes verification outputs to a temporary
directory:
\begin{lstlisting}
python3 -B reproduce.py
\end{lstlisting}
It verifies release hashes, reconstructs both literal tables, checks the
native/dense map, runs both finite fold checkers, reruns the comparison,
checks exact startup, and freshly derives all 5,084 runtime contracts.
The runtime phase took about 151 seconds in the preparation environment;
that elapsed time is environment-dependent, not a property of the machine.

\begin{lstlisting}
python3 -B reproduce.py --quick
\end{lstlisting}
The quick option explicitly skips the full symbolic runtime derivation.
It still checks the complete new finite fold theorem and startup, but is
not a substitute for the full replay when validating the primitive engine.
The local checker is additionally exercised with Python optimization enabled;
the inherited symbolic runtime refuses that mode because it uses assertions.

\begin{center}\small
\begin{tabularx}{\linewidth}{@{}p{34mm}X@{}}
\toprule
Path & Contents\\
\midrule
\path{report.tex}, \path{report.pdf} & This report; editable source and compiled output.\\
\path{paper/} & Generated complete dense transition table.\\
\path{machine/} & Final native/dense U208 tables, checked state map, unchanged source.\\
\path{baseline/} & Exact Joint U209 native parent and inherited runtime-label map.\\
\path{code/} & Deterministic builder, two finite checkers, comparison, startup,
and runtime replay.\\
\path{comparison/} & Exact table/source snapshots for all five code packets,
with their README and notices.\\
\path{evidence/} & Fresh receipts, full runtime cases, parent replay logs,
and comparison witnesses.\\
\path{reference/} & Unchanged Joint U209 ZIP with its complete lineage,
and the exact unary-theory ZIP.\\
\path{research/} & Rejected-rotation counterexample and bounded search scope.\\
\path{provenance/}, \path{licenses/} & Pinned input identities and retained attribution.\\
\path{SHA256SUMS} & Hashes of every other delivered file.\\
\bottomrule
\end{tabularx}
\end{center}

The PDF can be rebuilt with two runs of \texttt{pdflatex report.tex} in a
working copy. PDF bytes may differ with TeX versions and metadata; the
literal machine and source bytes are fixed.

\appendix
\section{Literal identities}\label{app:hashes}

\begin{description}[style=nextline]
\item[Parent Joint U209 native table]
\path{d0a155d3c683789f3bf5327c0cfb2a0cfd78076af8f1654d4ca426ccbb604d16}
\item[New U208 native table]
\path{527c2a6065ba7c9adb09294eed2b532810d812e46df4b5b912bae99ae6fb3a3a}
\item[New U208 dense table]
\path{9f4ea95849c5e1d01fa3f740023092230cd7006eed650b94a2d5da537e25212c}
\item[Unchanged source map]
\path{cddbed7f20e7a8b085b7ca0e76811d07b021522375643b5d1134eec0d2d23164}
\end{description}

\clearpage
\section{Complete dense U208 transition table}\label{app:table}

The working labels are $0,\ldots,207$; $H$ denotes the separate halt label
208. Each entry is write, move, target. This table is generated directly
from the delivered dense text table; it is not a schematic description.
Two state ranges are printed side by side.
\input{paper/transition_table.tex}

\begin{thebibliography}{10}\small
\bibitem{plain} ProveIt, \emph{U209 ConPA}, 29 October 2026.
\url{https://github.com/VladimirReshetnikov/ProveIt/blob/fe8b3a2d1db0033113e986590f79ecba3e5ef500/docs/incoming/U209_ConPA_2026-10-29.zip}.
\bibitem{composed} ProveIt, \emph{U209 ConPA Composed}, same snapshot.
\url{https://github.com/VladimirReshetnikov/ProveIt/blob/fe8b3a2d1db0033113e986590f79ecba3e5ef500/docs/incoming/U209_ConPA_Composed_2026-10-29.zip}.
\bibitem{folds} ProveIt, \emph{U209 ConPA Compatible Folds}, same snapshot.
\url{https://github.com/VladimirReshetnikov/ProveIt/blob/fe8b3a2d1db0033113e986590f79ecba3e5ef500/docs/incoming/U209_ConPA_Compatible_Folds_2026-10-29.zip}.
\bibitem{joint} ProveIt, \emph{U209 ConPA Joint}, same snapshot.
\url{https://github.com/VladimirReshetnikov/ProveIt/blob/fe8b3a2d1db0033113e986590f79ecba3e5ef500/docs/incoming/U209_ConPA_Joint_2026-10-29.zip}.
\bibitem{compat} ProveIt, \emph{U209 ConPA Compatibility}, same snapshot.
\url{https://github.com/VladimirReshetnikov/ProveIt/blob/fe8b3a2d1db0033113e986590f79ecba3e5ef500/docs/incoming/U209_ConPA_Compatibility_2026-10-29.zip}.
\bibitem{unary} ProveIt, \emph{Unary power covering at PA strength},
27 October 2026, archived at the same snapshot.
\url{https://github.com/VladimirReshetnikov/ProveIt/blob/fe8b3a2d1db0033113e986590f79ecba3e5ef500/docs/incoming/pa_unary_power_cover.zip}.
\bibitem{m254} Jeff's busy-beaver project, \emph{A verified 254-state
Turing machine that halts iff ZFC is inconsistent}, README, accessed
29 October 2026.
\url{https://github.com/jeff87654/busy-beaver/tree/zf-254/ZF/m254}.
\bibitem{k215} ProveIt, \emph{Formal correctness of the exact 215-state
machine}, pinned README.
\url{https://github.com/VladimirReshetnikov/ProveIt/tree/fe8b3a2d1db0033113e986590f79ecba3e5ef500/Computability/BusyBeaver/Research/ZFC/k225-formal/verification/k215}.
\bibitem{k215pa} Elliot Glazer, \emph{K215 halts iff ZFC is inconsistent
(Metamath Zero, PA)}, README, accessed 29 October 2026.
\url{https://github.com/elliotglazer/k215-pa}.
\end{thebibliography}
\end{document}
